\section{Disponibilização e Licenciamento}

Nestá seção descreve os mecanismos de disponibilização do dataset, incluindo o licenciamento adotado, as formas de acesso aos dados e a documentação associada. O objetivo é garantir que os usuários possam utilizar o dataset de maneira ética, legal e eficiente, promovendo a transparência e a colaboração na comunidade de pesquisa em segurança cibernética.

\subsection{Publicação}

Com o objetivo de garantir a reprodutibilidade dos resultados e facilitar a reutilização dos dados pela comunidade científica, o dataset desenvolvido neste trabalho foi disponibilizado publicamente por meio da plataforma
Zenodo. O Zenodo é um repositório de acesso aberto voltado ao armazenamento e preservação de produtos de pesquisa, como conjuntos de dados, softwares e publicações, atribuindo um Identificador de Objeto Digital (DOI) permanente a cada versão publicada.

A versão utilizada neste trabalho está disponível em:

\begin{itemize}
    \item DOI: \url{https://doi.org/10.5281/zenodo.21133620}
    \item Página do registro: \url{https://zenodo.org/records/21133620}
\end{itemize}

O repositório reúne o dataset final utilizado na construção do grafo de conhecimento, juntamente com um arquivo README com o dicionário de dados do dataset.

\subsection{Código-Fonte e Documentação}

Além do dataset, todos os scripts utilizados nas etapas de extração, transformação e integração dos dados foram disponibilizados publicamente em um repositório no GitHub. O repositório também contém a documentação do projeto, incluindo um arquivo \textit{README} com instruções para instalação, configuração do ambiente, execução do pipeline de processamento e importação dos dados para o Neo4j.

O código-fonte está disponível em:
\url{https://github.com/cybersecurity-datalake/cybersecurity-datalake}

\subsection{Licenciamento}

Todo o conteúdo disponibilizado neste trabalho é distribuído sob a licença MIT, uma licença permissiva de código aberto que permite o uso, cópia, modificação, distribuição e reutilização do material, inclusive para fins comerciais, desde que a declaração de direitos autorais e a licença original sejam preservadas.
